%! Orthogonality of the Four Subspaces — Unit 4, Lesson 4.1 — Slides (Beamer)
\documentclass[aspectratio=169, 11pt]{beamer}
\usepackage{linalg-beamer}   % provides \burgundyheader{} and \sectionlabel[color]{}
\usetikzlibrary{calc}
\newcommand{\vv}{\mathbf{v}}
\newcommand{\ww}{\mathbf{w}}
\newcommand{\xx}{\mathbf{x}}
\newcommand{\yy}{\mathbf{y}}
\newcommand{\zero}{\mathbf{0}}
\newcommand{\T}{\mathsf{T}}

\begin{document}

% TITLE SLIDE (CourseName is not defined in beamer, so the name is literal)
{
\setbeamercolor{background canvas}{bg=burgundy}
\begin{frame}[plain]
  \vspace{2.8cm}\hspace{0.55cm}%
  \begin{minipage}{0.9\textwidth}
    {\color{goldacc}\bfseries\small LESSON 4.1}\\[0.5cm]
    {\color{white}\bfseries\LARGE Orthogonality of the Four Subspaces}\\[1.0cm]
    {\color{blushmid}\small Linear Algebra~$\cdot$~Unit 4: Orthogonality}
  \end{minipage}
\end{frame}
}

% HOOK
\begin{frame}[t, plain]
  \burgundyheader{Which directions are perpendicular to the whole tabletop?}
  \sectionlabel{Hook}
  \begin{columns}[T]
    \begin{column}{0.56\textwidth}
      A flat tabletop is a \textbf{plane} in $\mathbb{R}^3$.\\[6pt]
      Only \emph{one} direction --- straight up/down --- is perpendicular to the \textbf{entire} plane:
      a single \textbf{line}.\\[8pt]
      Plane (dim $2$) $+$ line (dim $1$) $=$ the whole room ($\mathbf{2+1=3}$). Perpendicular \emph{and}
      filling: \textbf{orthogonal complements}.
    \end{column}
    \begin{column}{0.42\textwidth}
      \centering
      \begin{tikzpicture}[scale=0.9]
        % tabletop plane
        \draw[fill=blush, draw=burgundy, thick] (-1.3,-0.5)--(1.7,-0.5)--(1.3,0.5)--(-1.7,0.5)--cycle;
        \node[burgundy, font=\scriptsize] at (1.9,0.15) {plane};
        % vertical normal line
        \draw[->, very thick, royalblue] (0,0)--(0,2.0) node[right, font=\scriptsize]{$\perp$ line};
        \draw[->, very thick, royalblue] (0,0)--(0,-1.4);
        % right-angle mark
        \draw[charcoal, thin] (0,0.32)--(0.32,0.32)--(0.32,0);
      \end{tikzpicture}
    \end{column}
  \end{columns}
  \vspace{4pt}
  \begin{block}{The question of Lesson 4.1}
    Every matrix splits its input space into a row-space plane and a nullspace line just like this ---
    \textbf{why perpendicular, and why do they fill the space?}
  \end{block}
\end{frame}

% ORTHOGONAL SUBSPACES
\begin{frame}[t, plain]
  \burgundyheader{From orthogonal vectors to orthogonal subspaces}
  \sectionlabel{Definition}
  \begin{itemize}
    \setlength{\itemsep}{6pt}
    \item \textbf{Vectors:} $\vv\perp\ww$ means $\vv\cdot\ww=0$ (Lesson 1.2).
    \item \textbf{Subspaces:} $V\perp W$ means \emph{every} vector in $V$ is $\perp$ \emph{every} vector in $W$.
    \item \textbf{Shortcut:} enough to check each \emph{basis} vector of $V$ against each basis vector of $W$.
  \end{itemize}
  \vspace{4pt}
  \begin{block}{Caution --- perpendicular-looking is not enough}
    Two \textbf{walls} share a vertical line; a vector along it lies in both and can't be $\perp$ itself.
    So walls are \emph{not} orthogonal. Orthogonal subspaces meet \textbf{only at $\zero$}.
  \end{block}
\end{frame}

% WHY ROW SPACE PERP NULLSPACE
\begin{frame}[t, plain]
  \burgundyheader{Why the row space $\perp$ the nullspace}
  \sectionlabel[goldacc]{The ``why'' behind §3.6}
  Take $\xx$ in the nullspace: $A\xx=\zero$. Read it \textbf{one row at a time:}
  \[
    (\text{row }1)\cdot\xx=0,\quad (\text{row }2)\cdot\xx=0,\quad\dots,\quad (\text{row }m)\cdot\xx=0.
  \]
  \begin{itemize}
    \setlength{\itemsep}{5pt}
    \item So $\xx\perp$ \emph{every row} $\Rightarrow$ $\xx\perp$ every combination $\Rightarrow$ $\xx\perp$ the whole row space.
    \item $\boxed{\,N(A)\perp C(A^{\T})\,}$ inside $\mathbb{R}^n$.
    \item Same argument on $A^{\T}$ (rows of $A^{\T}=$ columns of $A$): $\boxed{\,N(A^{\T})\perp C(A)\,}$ inside $\mathbb{R}^m$.
  \end{itemize}
  \vspace{2pt}
  \begin{block}{One idea}
    The four-subspace pairing is just $\text{dot}=0$, repeated one row at a time.
  \end{block}
\end{frame}

% ORTHOGONAL COMPLEMENTS + BIG PICTURE
\begin{frame}[t, plain]
  \burgundyheader{Orthogonal complements: perpendicular \emph{and} filling}
  \sectionlabel{The big picture}
  \begin{columns}[T]
    \begin{column}{0.40\textwidth}
      In $\mathbb{R}^n$: \; $r+(n-r)=n$.\\[4pt]
      The two perpendicular pieces \textbf{fill} the space, so the nullspace is \emph{every} vector
      $\perp$ the row space --- its \textbf{orthogonal complement}.\\[6pt]
      \emph{Not} enough to be perpendicular: two $\perp$ lines in $\mathbb{R}^3$ give $1+1\neq3$.
    \end{column}
    \begin{column}{0.58\textwidth}
      \centering
      \begin{tikzpicture}[scale=0.72, every node/.style={font=\scriptsize}]
        \draw[burgundy, thick, rounded corners] (-0.3,-0.3) rectangle (3.9,4.6);
        \node[burgundy] at (1.8,4.95) {\textbf{Input } $\mathbb{R}^n$};
        \node[draw=burgundy, fill=blush, rounded corners, align=center, minimum width=3.0cm] at (1.8,3.5) {$C(A^{\T})$, dim $r$};
        \node[draw=burgundy, fill=blush, rounded corners, align=center, minimum width=3.0cm] at (1.8,1.0) {$N(A)$, dim $n{-}r$};
        \node[text=burgundy] at (1.8,2.25) {$\perp$};
        \draw[royalblue, thick, rounded corners] (5.9,-0.3) rectangle (10.1,4.6);
        \node[royalblue] at (8.0,4.95) {\textbf{Output } $\mathbb{R}^m$};
        \node[draw=royalblue, fill=skyblue, rounded corners, align=center, minimum width=3.0cm] at (8.0,3.5) {$C(A)$, dim $r$};
        \node[draw=royalblue, fill=skyblue, rounded corners, align=center, minimum width=3.0cm] at (8.0,1.0) {$N(A^{\T})$, dim $m{-}r$};
        \node[text=royalblue] at (8.0,2.25) {$\perp$};
        \draw[->, very thick, charcoal] (4.1,2.15)--(5.7,2.15) node[midway, above]{$A$};
      \end{tikzpicture}
    \end{column}
  \end{columns}
\end{frame}

% WORKED EXAMPLE
\begin{frame}[t, plain]
  \burgundyheader{One matrix, all four subspaces}
  \sectionlabel[goldacc]{Check it}
  $A=\begin{bsmallmatrix} 1 & 1 & 2\\ 2 & 1 & 3\\ 3 & 2 & 5 \end{bsmallmatrix}$, rank $r=2$, with
  nullspace direction $\xx=\begin{bsmallmatrix} -1\\-1\\1 \end{bsmallmatrix}$ and left-nullspace
  direction $\yy=\begin{bsmallmatrix} 1\\1\\-1 \end{bsmallmatrix}$.
  \begin{columns}[T]
    \begin{column}{0.5\textwidth}
      \begin{block}{Input room ($\mathbb{R}^3$)}
        row$_1\cdot\xx=\begin{bsmallmatrix} 1\\1\\2 \end{bsmallmatrix}\cdot\begin{bsmallmatrix} -1\\-1\\1 \end{bsmallmatrix}=-1-1+2=0.$\\[4pt]
        Plane (dim $2$) $\perp$ line (dim $1$), $2+1=3$.
      \end{block}
    \end{column}
    \begin{column}{0.5\textwidth}
      \begin{block}{Output room ($\mathbb{R}^3$)}
        col$_1\cdot\yy=\begin{bsmallmatrix} 1\\2\\3 \end{bsmallmatrix}\cdot\begin{bsmallmatrix} 1\\1\\-1 \end{bsmallmatrix}=1+2-3=0.$\\[4pt]
        Plane (dim $2$) $\perp$ line (dim $1$), $2+1=3$.
      \end{block}
    \end{column}
  \end{columns}
  \vspace{4pt}
  \begin{block}{Next --- Lesson 4.2}
    The error of a best guess lands in the \emph{complement}. We'll \emph{use} that: force the error
    $\perp$ the subspace (dot $=0$) to find the closest point --- the \textbf{projection}.
  \end{block}
\end{frame}

\end{document}
